% Phase 4: Empirical Evaluation and Comparative Analysis
% Author: Abdullah Zia (Evaluation, Analysis, and Report Lead)

\section{Phase 4: Empirical Evaluation}

\subsection{Evaluation Methodology and Metric Formulation}
To rigorously assess the performance of the neural-guided heuristic relative to classical
baselines, we establish an evaluation protocol spanning regression fidelity, auxiliary
classification discrimination, and end-to-end graph search efficiency.

\subsubsection{Data Leakage Prevention and Split Structure}
Evaluation is strictly isolated on the unseen test set of 282 mazes partition (15\% of
the cleaned corpus), partitioned by complete maze environment rather than individual cell
instances. This environment-level partitioning guarantees zero leakage of global maze
topologies, wall distributions, and goal placements between training and evaluation phases.

\subsubsection{Regression and Admissibility Metrics}
For a given maze of dimension $N \times N$, let $h^\ast(n)$ denote the ground-truth optimal
shortest-path cost from cell $n$ to the goal, and let $\hat{h}_\theta(n)$ denote the
neural prediction. Because models are trained to predict normalized distances $y = h^\ast / N$,
primary regression metrics include:
\begin{itemize}
    \item \textbf{Mean Squared Error (MSE)} and \textbf{Root Mean Squared Error (RMSE)}:
    \begin{equation}
        \text{RMSE} = \sqrt{\frac{1}{|\mathcal{S}|} \sum_{n \in \mathcal{S}} \left(\hat{y}(n) - y(n)\right)^2}
    \end{equation}
    where $\mathcal{S}$ is the set of all free reachable cells.
    \item \textbf{Mean Absolute Error (MAE)} in normalized units and raw grid cells ($\text{MAE} \times N$).
    \item \textbf{Correlation Measures}: Pearson linear correlation coefficient $r$ and Spearman rank correlation coefficient $\rho$.
    \item \textbf{Admissibility Violation Rate}: Fraction of reachable cells where the unscaled heuristic overestimates true distance beyond tolerance $\epsilon$:
    \begin{equation}
        \text{Violations}(\epsilon) = \frac{1}{|\mathcal{S}|} \sum_{n \in \mathcal{S}} \mathbb{I}\left[(\hat{h}_\theta(n) - h^\ast(n)) > \epsilon\right]
    \end{equation}
\end{itemize}

\subsubsection{Auxiliary 10-Bin Classification Metrics}
To assess spatial range discrimination, targets are discretized into 10 decile distance bins.
We compute per-class Precision, Recall, and $F_1$-score, macro-averaged and weighted $F_1$,
and Macro One-vs-Rest Area Under the ROC Curve (AUC-ROC).

\subsubsection{Search Performance Metrics}
During deterministic $A^\ast$ execution, we measure:
\begin{enumerate}
    \item \textbf{Nodes Expanded}: The exact count of states popped from the priority queue and closed.
    \item \textbf{Nodes Generated}: Total valid states inserted into the open set.
    \item \textbf{Path Cost and Optimality Ratio}: Ratio of resulting path length to optimal Dijkstra cost $c / c^\ast$.
    \item \textbf{Relative Node Reduction}: Reduction in expansions relative to Manhattan $A^\ast$:
    \begin{equation}
        \Delta_{\text{expanded}} = 1 - \frac{\text{Nodes}_{\text{neural}}}{\text{Nodes}_{\text{manhattan}}}
    \end{equation}
    \item \textbf{Runtime and Peak Open Set}: CPU wall-clock latency per maze and memory footprint.
\end{enumerate}
